\chapter{结论与后续工作}

\section{主要贡献}

本研究以"KAN符号回归中物理气象变量为何被系统性剪除"为核心问题，通过递进式实验设计，在机制揭示、诊断方法和解决路径三个层面取得了以下贡献。

\textbf{贡献一：建立了从KAN教师到符号公式的完整提取管线并验证了KAN的预测有效性。}在ARPA-E PERFORM数据集的净负荷差分预测任务上，当前canonical matched主实验设置中的KAN教师模型达到了skill=0.453的预测技能分数，在统计显著性检验中优于参数对齐的MLP基线（skill=0.430，Wilcoxon检验$p=0.0005$），同时大幅超过特征对齐PySR（skill=0.076）和持续性基线。

\textbf{贡献二：发现并通过干预实验支持了自回归捷径竞争是破坏符号可辨识性的主导机制。}在直接净负荷目标上，9组符号配置的提取结果全部坍缩为仅含负荷项的极简公式（$0.0007 \times \text{load} - 3.0$），VER=0/9。S2自回归阻断实验提供了分层证据：Case 3给出强干预支持（$\Delta\VER=+0.60$，95\% CI=[0.20, 1.00]）；严格matched的Case 4 direct-task对照仅在4个有效paired seed上得到$\Delta\VER=+0.25$（95\% CI=[0.00, 0.75]），另有1个seed两次未产出post-prune checkpoint，说明这一路证据更弱且不稳定。

\textbf{贡献三：提出了面向符号可辨识性的VER/FAR诊断指标体系。}传统符号回归评价指标（$R^2$、复杂度）关注公式的拟合质量和简洁度，但不考察特定物理变量是否出现在公式中。本研究定义的VER和FAR从稀疏结构和符号公式两个层面量化了物理变量的可辨识性，使"可辨识性"从定性描述转化为可度量、可比较的定量判据。

\textbf{贡献四：设计了S3结构化分解方案，并完成了净负荷总公式的显式闭环复算。}S3结构化组合预测器的预测技能分数（skill=0.515）已经优于直接KAN教师（skill=0.453），RMSE为1254.62；同时三个子任务的FAR均达到3/3。更重要的是，本研究已经显式构造$F_{\text{net}} = F_{\text{load}} - F_{\text{wind}} - F_{\text{solar}}$并完成test集逐点复算，得到S3总公式的RMSE为2644.36、skill为-0.021。由此，本文不再停留在"三条局部公式 + 组合预测器"的半闭环状态，而是进一步揭示了公式对象相对预测器对象存在明显transfer gap这一新的事实边界。

\textbf{贡献五：诚实记录了方法的边界失败案例。}太阳辐照方面，VER(GHI)=3/3的正面结果主要成立于部分短中期设置；风电方面，wind更适合作为边界案例而非"训练失败"叙事，其可辨识性随预测时域呈现非单调变化；其他物理量在长horizon下只宜表述为整体更不稳定、部分减弱；S2直接阻断虽可恢复可辨识性，但伴随明显精度代价。

\section{收口工作与后续研究}

\textbf{S3总公式的工程收口已经完成，当前剩余的问题不再是"公式是否存在"，而是"公式相对预测器为何仍有明显transfer gap"。}本文已经完成
\begin{equation}
F_{\text{net}} = F_{\text{load}} - F_{\text{wind}} - F_{\text{solar}}
\end{equation}
的显式合成与test集逐点复算，因此现在可以把S3称为"已闭环的净负荷总公式方案"。但复算结果也显示，总公式的RMSE（2644.36）显著高于结构化组合预测器（1254.62）。因此，后续研究的首要任务不再是补一条总公式，而是解释并缩小这一从"局部公式有效"到"全局总公式数值可用"之间的迁移间隙。

\textbf{在补充交付之外，后续研究仍可沿三个方向推进。}其一，重构严格时间因果的数据管线，将插值放在训练/验证/测试划分之后，以进一步收紧机制结论的边界。其二，在更多区域、更多时间分辨率的数据集上复现S2与S3，以检验自回归捷径竞争及wind边界现象的普遍性。其三，在KAN训练中引入物理约束或更丰富的符号基元，探索是否可以在不牺牲精度的前提下进一步降低局部公式的迁移间隙和复杂度。
